\begin{question}
\noindent The diagram shows triangle $T$ and its images $A$, $B$ and $C$ after three different single transformations.

\begin{center}
\begin{tikzpicture}[scale=0.62, transform shape]
    \def\axisLimit{8}

    \definecolor{borderColour}{HTML}{000000}
    \definecolor{fillColourT}{HTML}{BFD8FF}
    \definecolor{fillColourA}{HTML}{FFD8B0}
    \definecolor{fillColourB}{HTML}{D8D8D8}
    \definecolor{fillColourC}{HTML}{C8E6C9}

    % Draw grid
    \draw[step=1cm,gray,very thin] (-\axisLimit,-\axisLimit) grid (\axisLimit,\axisLimit);

    % Draw axes
    \draw[thick,->] (-\axisLimit-0.5,0) -- (\axisLimit+0.5,0) node[anchor=north] {$x$};
    \draw[thick,->] (0,-\axisLimit-0.5) -- (0,\axisLimit+0.5) node[anchor=east] {$y$};

    % Label x-axis and y-axis
    \foreach \x in {-\axisLimit,...,\axisLimit}
        \draw (\x cm,1pt) -- (\x cm,-1pt) node[anchor=north] {$\x$};
    \foreach \y in {-\axisLimit,...,\axisLimit}
        \draw (1pt,\y cm) -- (-1pt,\y cm) node[anchor=east] {$\y$};

    % Triangle T: P(1,1), Q(1,3), R(3,1)
    \coordinate (P) at (1,1);
    \coordinate (Q) at (1,3);
    \coordinate (R) at (3,1);
    \filldraw[fill=fillColourT, draw=borderColour, line width=1pt] (P) -- (Q) -- (R) -- cycle;
    \node[below left] at (P) {$P$};
    \node[above left] at (Q) {$Q$};
    \node[below right] at (R) {$R$};
    \node at (1.9,1.8) {$T$};

    % Triangle A: translation of T by (-6,-6): P1(-5,-5), Q1(-5,-3), R1(-3,-5)
    \coordinate (P1) at (-5,-5);
    \coordinate (Q1) at (-5,-3);
    \coordinate (R1) at (-3,-5);
    \filldraw[fill=fillColourA, draw=borderColour, line width=1pt] (P1) -- (Q1) -- (R1) -- cycle;
    \node[below left] at (P1) {$P_1$};
    \node[above left] at (Q1) {$Q_1$};
    \node[below right] at (R1) {$R_1$};
    \node at (-4.1,-4.2) {$A$};

    % Triangle B: reflection of T in y-axis: P2(-1,1), Q2(-1,3), R2(-3,1)
    \coordinate (P2) at (-1,1);
    \coordinate (Q2) at (-1,3);
    \coordinate (R2) at (-3,1);
    \filldraw[fill=fillColourB, draw=borderColour, line width=1pt] (P2) -- (Q2) -- (R2) -- cycle;
    \node[below right] at (P2) {$P_2$};
    \node[above right] at (Q2) {$Q_2$};
    \node[below left] at (R2) {$R_2$};
    \node at (-1.9,1.8) {$B$};

    % Triangle C: enlargement of T, scale factor 2, centre (5,-1)
    % P(1,1)->(5,-1)+2*((1,1)-(5,-1))=(5,-1)+2*(-4,2)=(5-8,-1+4)=(-3,3)... need within grid, recompute centre
    % Let's choose centre (5,1): image of P = (5,1)+2*((1,1)-(5,1)) = (5,1)+2*(-4,0)=(5-8,1)=(-3,1) -- overlaps B, avoid
    % Choose centre (7,-3): P'=(7,-3)+2*((1,1)-(7,-3))=(7,-3)+2*(-6,4)=(7-12,-3+8)=(-5,5) out of nice zone
    % Simpler: centre (5,-3), scale factor 2
    % P(1,1): P' = (5,-3)+2*((1,1)-(5,-3)) = (5,-3)+2*(-4,4) = (5-8,-3+8) = (-3,5)
    % Q(1,3): Q' = (5,-3)+2*((1,3)-(5,-3)) = (5,-3)+2*(-4,6) = (5-8,-3+12) = (-3,9) -- too high
    % Use centre (6,4), scale factor 2 (placed top right, away from others)
    \coordinate (centreC) at (6,4);
    \coordinate (P3) at ($(centreC)+2*(P)-2*(centreC)$);
    \coordinate (Q3) at ($(centreC)+2*(Q)-2*(centreC)$);
    \coordinate (R3) at ($(centreC)+2*(R)-2*(centreC)$);
    \filldraw[fill=fillColourC, draw=borderColour, line width=1pt] (P3) -- (Q3) -- (R3) -- cycle;
    \node[below left] at (P3) {$P_3$};
    \node[above left] at (Q3) {$Q_3$};
    \node[below right] at (R3) {$R_3$};
    \node at (0,4.8) {$C$};

\end{tikzpicture}
\end{center}

\noindent Triangle $T$ has vertices $P(1,1)$, $Q(1,3)$ and $R(3,1)$.

\begin{enumerate}
    \item[(a)] Triangle $A$ (vertices $P_1, Q_1, R_1$) is the image of $T$ under a single transformation. Describe this transformation fully.

    \vspace{3cm}
    \answerlines{2}{2}

    \item[(b)] Triangle $B$ (vertices $P_2, Q_2, R_2$) is the image of $T$ under a single transformation. Describe this transformation fully.

    \vspace{3cm}
    \answerlines{2}{2}

    \item[(c)] Triangle $C$ (vertices $P_3, Q_3, R_3$) is the image of $T$ under a single transformation. Describe this transformation fully, stating the coordinates of the centre.

    \vspace{3cm}
    \answerlines{2}{1}
\end{enumerate}
\end{question}
